%%%%%%%%%%%%%%%%%%%%%%%%%%%%%%%%%%%%%%%%%
% "ModernCV" CV and Cover Letter
% LaTeX Template
% Version 1.11 (19/6/14)
%
% This template has been downloaded from:
% http://www.LaTeXTemplates.com
%
% Original author:
% Xavier Danaux (xdanaux@gmail.com)
%
% License:
% CC BY-NC-SA 3.0 (http://creativecommons.org/licenses/by-nc-sa/3.0/)
%
% Important note:
% This template requires the moderncv.cls and .sty files to be in the same 
% directory as this .tex file. These files provide the resume style and themes 
% used for structuring the document.
%
%%%%%%%%%%%%%%%%%%%%%%%%%%%%%%%%%%%%%%%%%

%----------------------------------------------------------------------------------------
%	PACKAGES AND OTHER DOCUMENT CONFIGURATIONS
%----------------------------------------------------------------------------------------

\documentclass[11pt,a4paper,sans]{moderncv} % Font sizes: 10, 11, or 12; paper sizes: a4paper, letterpaper, a5paper, legalpaper, executivepaper or landscape; font families: sans or roman

\moderncvstyle{casual} % CV theme - options include: 'casual' (default), 'classic', 'oldstyle' and 'banking'
\moderncvcolor{blue} % CV color - options include: 'blue' (default), 'orange', 'green', 'red', 'purple', 'grey' and 'black'

\usepackage{lipsum} % Used for inserting dummy 'Lorem ipsum' text into the template

\usepackage[scale=0.75]{geometry} % Reduce document margins
%\setlength{\hintscolumnwidth}{3cm} % Uncomment to change the width of the dates column
%\setlength{\makecvtitlenamewidth}{10cm} % For the 'classic' style, uncomment to adjust the width of the space allocated to your name

%----------------------------------------------------------------------------------------
%	NAME AND CONTACT INFORMATION SECTION
%----------------------------------------------------------------------------------------

\firstname{Arseniy} % Your first name
\familyname{Trushin} % Your last name

% All information in this block is optional, comment out any lines you don't need
\title{Curriculum Vitae}
%\address{m. Sukharevskaya pl, 3, 128}{Moscow, Russia 129090}
\mobile{+7(916) 588 5156}
% \phone{+7(495) 694 1520}
%\fax{(000) 111 1113}
\email{tarstars@gmail.com}
\homepage{www.linkedin.com/in/tarstars}{www.linkedin.com/in/tarstars}
%\extrainfo{additional information}
\photo[70pt][0.4pt]{pictures/trushin} % The first bracket is the picture 
%height, the second is the thickness of the frame around the picture (0pt for 
%no frame)
%\quote{"A witty and playful quotation" - John Smith}

%----------------------------------------------------------------------------------------

\begin{document}

\makecvtitle % Print the CV title

%----------------------------------------------------------------------------------------
%	EDUCATION SECTION
%----------------------------------------------------------------------------------------

%todo: 
%*) В списке проектов структурировать про проектам. Для каждого указать, 
%   в чём были челленджи и достижения.
%*) Про анализ стихотворений: analyst + software development
%*) line count - убрать
%*) про систему навигации по звёздам - для микроконтроллера, жёсткие 
%       требования, на спутник 
%*) учебная деятельность - перенести повыше, записать как проект
%*) выкинуть про книжки
%*) topcoder - в interests
%*) скилы - в один список ключевых слов
%*) operating systems

\section{Academic Qualifications}
%\textit{GPA -- 4.5}
\cventry{2001}{Specialist's degree}{Degree in physics, The Physical 
Faculty of Lomonosov Moscow State University}{Moscow}{}{In specialty 
``Physics''}
\cventry{2004}{Ph.D.}{ ``Yb photoluminescence in ZnSe-ZnCdSe films''}{The 
Physical Faculty of Lomonosov Moscow State University}{Moscow}{In specialty ``Laser Physics''}

%----------------------------------------------------------------------------------------
%	WORK EXPERIENCE SECTION
%----------------------------------------------------------------------------------------

\section{Experience}

\subsection{Vocational}

\cventry{August 2025 -- now}{Lead Developer}{Yandex Delivery}{Agentic system for support operators}{}{}

\cventry{November 2022 -- August 2025}{Lead Developer}{Yandex Schoolbook}{}{}{}

\cventry{July 2019 --  November 2022}{Senior Programmer}{App in the air}{}{}{}

\cventry{March 2019 -- June 2019}{Lead Programmer}{Joom}{Search department}{}{}

\cventry{March 2018 -- March 2019}{Data Scientist}{HeadHunter}{Search department}{}{}

\cventry{January 2018 -- March 2018}{Lead Analyst}{Equifax}{Credit fraud prevention}{}{}

\cventry{May 2016 -- January 2018}{Programmer}{Yandex}{Search factors development. Search engine support}{}{}

\cventry{May 2015 -- May 2016}{Programmer}{Moscow office of Cadence design 
systems}{Development of complex software for physical simulation}{}{}

\cventry{May 2009 -- Sep 2014}{Assistant professor}{Physical Department of 
MSU}{Modeling of acoustical and optical 
processes in crystals}{}{}

\section{Teaching and Supervision}

\cvitem{2025}{Taught a course on Machine Learning at Russian-Armenian University (RAU).}

\cvitem{2025}{Supervised four graduate theses:
  \begin{itemize}
    \item \textit{Uncompromised efficiency of the linear heuristics} (Russian-Armenian University)
    \item \textit{SOTA for airline traffic in the USA} (Russian-Armenian University)
    \item \textit{How to train your LLM to play QUIXO} (Moscow Institute of Physics and Technology)
    \item \textit{Teaching LLM to play QUIXO} (Moscow Institute of Physics and Technology)
  \end{itemize}
}

\subsection{Programming projects}

\cventry{August 2025 -- now}{Lead Developer}{Yandex Delivery}{Agentic system for support operators}{}{
 Working on an agentic system that helps support operators use internal manuals and resolve user requests. \textbf{Achievements}
 \begin{itemize}
    \item \textbf{Product research and specification}: Translated an ambiguous product request into a clear, actionable task specification through product research and scenario analysis.
    \item \textbf{LLM-powered manual conversion}: Converted 150 support manuals into graph-based descriptions using LLM-powered agents, making manual knowledge easier to structure, validate, and use in agentic workflows.
    \item \textbf{Workflow acceleration}: Delivered a 2x speedup in the manual conversion workflow.
    \item \textbf{Production application}: Implemented the approach as a production-grade web-based agentic application and supported colleagues in integrating it into their workflow.
 \end{itemize}
}

\cventry{November 2022 -- August 2025}{Lead Developer}{Yandex Schoolbook}{}{}{
 Main responsibilities include training and evaluating large language models (LLMs) for educational applications, leading projects, and enhancing existing systems. \textbf{Achievements}
 \begin{itemize}
    \item \textbf{Metrics for R\&D Projects}:  Designed and led a 3-month R\&D project involving 7 contributors and 25+ tickets. The project aimed to correlate user satisfaction metrics (likes/dislikes) with offline quality benchmarks for three LLM use cases (\textit{explain terms}, \textit{debug code}, \textit{dialogs}, 150k MAU total, $\sim$3k MAU per scenario).  Found that a simple filter in \textit{explain terms} boosted likes/dislikes by +9\%; model retraining in \textit{debug code} improved offline quality from 68\% to 87\% (+19\%) and also raised likes/dislikes by +9\%. Identified key pain points in \textit{dialogs} (identity confusion, cyclic behavior) with fixes projected to improve feedback by +7\%.
    \item \textbf{Processed and implemented a product request}: triaged user scenarios.
    \item \textbf{LLM Assessment and Reassesment for Unified State Exams}: Created a basket of 105 tasks and a method of model's evaluation. Conducted an evaluation for an internal LLM and repeated it one year later for an updated version. Results of these evaluations was used to estimate viability of department projects. Led a team of two members.
    \item \textbf{Math Tutor Technology Rollout for EGE Platform}: Led the deployment of a new dialog-based tutoring technology for the Unified State Exam platform. Coordinated task ingestion, software development, offline evaluation, and online A/B testing. Resulted in +21\% gain (72\% $\rightarrow$ 93\%) on offline Hanmigo-style dialog metrics; dialogs in the platform now run on the new infrastructure.
    \item \textbf{Created a fallback model for production}: prepared datasets (auto/manual annotation), trained models, ran quality evaluations, and prepared deployment materials.
    \item \textbf{Developed a synthetic dialogue generator} to create artificial dialog datasets for LLM pretraining and evaluation.
    \item \textbf{Dialog estimation application} Python + QT application for chatting with LLM and evaluating the results of the chat. It's the first version of the program the team has been using for two years.
    \item \textbf{Organized ML Artifacts Storage System}: Implemented a centralized system for storing machine learning artifacts within the team, improving collaboration and efficiency in model development and deployment.
    \item \textbf{EGE Debugging Initiative}: Proposed a methodology for collecting evaluation and training datasets for LLMs related to the Unified State Exam (EGE). Established processes for result evaluation and introduced a new quality metric to improve model assessment.
    \item \textbf{Skill Graph Refactoring}: Analyzed and refactored a legacy algorithm for the skill graph module. Rewrote the system using modern technologies and integrated it with a new data source, enhancing performance and maintainability. Data freshness: 20 hours -> 2 hours.
    \item \textbf{'Success Rate' Analytical Project}: Conducted an analytical project to compute 'Success Rate' metrics by analyzing logs of student activity. Tested the hypothesis that the number of successfully solved tasks differs between mobile and desktop platforms, providing insights for platform optimization.
 \end{itemize}
}

\cventry{July 2019 -  November 2022}{Analytics Department}{App In The Air}{Simulation of users behaviour.
\textbf{Challenges} Improve performance and code quality of the Python program. 
\textbf{Achievements} I managed an ``Automatic Insights for airlines'' project from scratch. I implemented the Extrapolating Gradient Boosted Decision Trees algorithm in Golang and applied it for company's data. I proposed several architectural solutions in the user statistics report project. I moved a code of users behaviour simulator from jupyter notebooks to .py files and covered it with unit tests}{}{}


\cventry{March 2019 - June 2019}{Factors development}{Joom}{Implementation of algorithms for factors calculations.
\textbf{Challenges} New languages - Scala and Golang. Fuzzy task on ngram query split. 
\textbf{Achievements} Scala program for data download were completed as well as Golang part of realtime calculations. Two algorithms proposed: dynamic programming and Byes tokenization
\textbf{Leave} Due to illness during the trial period}{}{}

\cventry{April 2018 - May 2018}{User clicks prediction}{Pet project}{Tensorflow model was trained for prediction of user click probability distribution for the given screenshot of a browser.
\textbf{Challenges} Tensorflow was absolutely unknown instrument.
\textbf{Achievements} Model was trained and evaluated. U-net with the elements of dense-net network architecture was reinvented}{}{}

\cventry{March 2018 - March 2019}{Search engine improvement}{Headhunter}{Implementation of algorithms for factors calculations.
\textbf{Challenges} Very interesting and clever but mostly undocumented python code.
\textbf{Achievements} Several enhancements of search engine were implemented.
\textbf{Leave} Job proposal with better conditions}{}{}

\cventry{January 2018 - March 2018}{Credit Fraud Prevention}{Equifax}{Oracle calculations
  and testing of factors for credit fraud detection. \textbf{Challenges} I have never written such complex sql queries before.
\textbf{Achievements} I learn how to write queries for factors calculations and built ML model for credit fraud prediction.
\textbf{Leave} My ML expertise was not valued at this company.}{}{}

\cventry{May 2016 - January 2018}{Factors development. Search core support}{Yandex}{Map-Reduce calculations
and testing of behaviorial factors. Support of links data. \textbf{Challenges} A lot to learn: python and
c++ interfaces for Yandex Map-Reduce technology, search logs formats and access libraries. Legacy C++
code for links database. \textbf{Achievements} Calculation of several models. Successfull upgrade of links database.
\textbf{Leave} Job proposal with better conditions}{}{}

\cventry{May 2015 - May 2016}{Randomized field solver}{Cadence}{Random value properties estimation.
Python script to run and monitore execution of big tasks on inhomegenious computational resources.
Analyst, software engineer. \textbf{Challenges}: ill-formulated problem for random value properties estimation,
gradually changing requirements for multi-tasking module. \textbf{Achievements}: 
most effective statistic methods for the hypotheses testing were found; python module for 
dependencies resolution and monitoring of tasks execution in nonhomogenious computational
facilities
\textbf{Leave} Job proposal with better conditions}{}{}

\cventry{May 2009 - June 2009}{Denoising of optical profilometer picture}{Scientific tool}{ An algorithm of
signal denoising by averaging in one direction was implemented in C++ and Golang
\textbf{Challenges} Strict time limitations on development. Lack of convenient versatile instruments like PIL in jupyter notebook
\textbf{Achievements} Low noise results was obtained from the images}{}{}

\cventry{May 2009 - May 2014}{Beam structure}{Physical Faculty}{Software for simulation 
of excitation and propagation of acoustical beams in crystalline media. 
Analyst, software engineer. \textbf{Challenges}: complex math scheme, 3d visualization 
of calculation results, resource-consuming tasks. \textbf{Achievements}: 
picturesque opengl-based interactive model of the slowness surface;
FFT-based scheme for an acoustic field calculation}{}{}

\cventry{2006 - 2009}{Navigation system}{Lebedev Physical Institute}{Star 
navigation system on satellite}{\textbf{Challenges}: strict resource (time and 
memory) limitations, \textbf{Achievements}: calculation scheme was translated 
from IDE to asm; python compiler for assembler was written}{}{}

\cventry{2003-2004}{Rhyme analysis}{}{Analyst, software 
engineer. \textbf{Challenges}: very fuzzy task specification} 
{\textbf{Achievements}: strict formulation 
of algorithms and their implementation}{}{}

\cventry{2001-2003}{X-ray analysis}{}{Analyst, software engineer} 
{\textbf{Challenges}: uncertain description of interaction between user and 
program, necessity of serialization of unfinished session}{\textbf{Achievements}: full 
specification of user interaction and GUI was devised and implemented; object 
model of data was constructed; file format for session serialization was 
proposed}{}

\cventry{1996--2001}{Student}{Physical Faculty of Moscow State 
University}{Diploma projects ``Diffraction of X-rays on crystals deformed by 
ion-implantation '', ``Alignment of retina-scans''} {\textbf{Challenges}: no 
strict task description} {\textbf{Achievements}: full task description was 
proposed, task was implemented} {}


%-------------------------------------------------------------------------------
%---------
%	AWARDS SECTION
%-------------------------------------------------------------------------------
---------
%\section{Awards}

%\cvitem{2004}{Fenia Berz Award}
%\cvitem{2010}{``The best presentation on the conference'' - Coherent Optics 
%and 
%Optical Spectroscopy}

%----------------------------------------------------------------------------------------
%	COMPUTER SKILLS SECTION
%----------------------------------------------------------------------------------------

\newpage

\section{Computer skills}

% win32api, mfc, html, R,
%
\cvitem{Keywords}{c++, golang, python, git, svn, perforce, stl, bash,  
ssh, qt, povray, opengl, scipy, pandas, numpy, sklearn}

%----------------------------------------------------------------------------------------
%	COMMUNICATION SKILLS SECTION
%----------------------------------------------------------------------------------------

\section{Communication Skills}

\cvitem{2000-2014}{Over 20 oral presentations at different conferences 
including 4 international conferences with English presentations}

%-------------------------------------------------------------------------------
%---------
%	LANGUAGES SECTION
%----------------------------------------------------------------------------------------

\section{Languages}

\cvitemwithcomment{Russian}{Mothertongue}{}
\cvitemwithcomment{English}{Upper Intermediate}{Reading fluent, Conversationally
fluent}

%----------------------------------------------------------------------------------------
%	INTERESTS SECTION
%----------------------------------------------------------------------------------------

\section{Miscellaneous}
\cvitem{Topcoder login}{142857}
\cvitem{Machine Learning}{Vorontsov's course in ``School of data analysis''. Express course of Mehryar Mohri in Yandex.}

\section{Interests}

\renewcommand{\listitemsymbol}{-~} % Changes the symbol used for lists

I like to ride my bike with planetary gearbox. I have canon 1000d camera and 
sometimes I use it ;)  

%----------------------------------------------------------------------------------------
%	COVER LETTER
%----------------------------------------------------------------------------------------

% To remove the cover letter, comment out this entire block

%\clearpage

%\recipient{HR Department}{Corporation\\123 Pleasant Lane\\12345 City, State} % 
%Letter recipient
%\date{\today} % Letter date
%\opening{Dear Sir or Madam,} % Opening greeting
%\closing{Sincerely yours,} % Closing phrase
%\enclosure[Attached]{curriculum vit\ae{}} % List of enclosed documents

%\makelettertitle % Print letter title

%\lipsum[1-3] % Dummy text

%\makeletterclosing % Print letter signature

%----------------------------------------------------------------------------------------

\end{document}
